%==============================================================================
%  CAPÍTULO XIII
%==============================================================================
\chapter{Derecho concursal internacional e insolvencia transfronteriza}
\label{cap:transfronteriza}

\fichaCapitulo{Cooperación judicial y extraterritorialidad.}{Adopción de la Ley Modelo
UNCITRAL, determinación del \comi{} y análisis de los casos LATAM y Astaldi.}

%==============================================================================
\bloqueTeorico
%==============================================================================

\subsection{De la territorialidad al universalismo cooperativo}
\pendiente{Desarrollar: superación del principio de jurisdicción exclusiva sobre los
activos locales.}

\subsection{La Ley Modelo UNCITRAL en el Capítulo VIII de la Ley \Nro 20.720}
La adopción del modelo se produjo mediante la incorporación de la Ley Modelo de la
CNUDMI \parencite{ley-modelo-uncitral} al Capítulo VIII de la \LIR{}. Para un balance de
su primera década de aplicación en Chile, véase \textcite{doc-diez-anos-uncitral}.
\pendiente{Desarrollar: herramientas frente al ocultamiento transfronterizo de bienes.}

%==============================================================================
\bloqueNormativo
%==============================================================================

\subsection{Principios rectores}

Cooperación judicial, seguridad jurídica para la inversión, administración equitativa del
concurso y protección de los activos de deudores multinacionales.

\subsection{La determinación del centro de intereses principales}

Presupuestos de visibilidad externa y establecimiento permanente para fijar la
jurisdicción rectora del concurso.

\subsection{Reconocimiento del procedimiento extranjero}

Distinción entre procedimiento extranjero principal y no principal, y efectos de cada
uno. Sobre la práctica chilena de reconocimiento, véase
\textcite{doc-reconocimiento-extranjero}; para una guía de tramitación,
\textcite{doc-transfronteriza-guia}.

\subsection{Medidas provisionales urgentes}

Suspensión de ejecuciones, prohibición de transferir activos locales y suministro de
información patrimonial, otorgables tras la sola presentación de la
solicitud.\verificar{Confirmar el artículo que regula las medidas y su catálogo.}

%==============================================================================
\bloquePractico
%==============================================================================

\subsection{Solicitud de reconocimiento ante el juzgado civil chileno}

\begin{enumerate}
  \item \textbf{Presentación de la solicitud.} Escrito ante el juzgado de letras en lo
  civil competente, acompañando copia certificada de la resolución extranjera que declara
  la insolvencia y acredita la designación del representante extranjero.

  \item \textbf{Declaración de concursos paralelos.} Declaración jurada que individualice
  todos los procedimientos de insolvencia iniciados respecto del deudor de que tenga
  conocimiento el representante extranjero.

  \item \textbf{Medidas provisionales.} Solicitud incidental de paralización de las
  ejecuciones individuales dirigidas contra los activos locales y, en su caso,
  designación de un interventor.

  \item \textbf{Publicidad y oposición de terceros.} Publicación en el \boletin{} y
  sustanciación de las controversias de los acreedores locales.
\end{enumerate}

\subsection{Casos paradigmáticos de la práctica chilena}

\begin{table}[htbp]
\centering
\small
\caption{Casos de referencia en materia de insolvencia transfronteriza}
\label{tab:cap13-casos}
\begin{tabularx}{\textwidth}{@{}>{\raggedright\arraybackslash}p{0.16\textwidth} XXX@{}}
\toprule
\textbf{Caso} & \textbf{Tribunal e hitos} & \textbf{Criterios aplicados}
& \textbf{Consecuencia práctica} \\
\midrule
Chang Rajii & 15.\textsuperscript{o} Juzgado Civil de Santiago (2016)
& Cooperación judicial internacional por vía diplomática para incautar activos bancarios
e inmobiliarios situados en el extranjero
& Extensión de las facultades del liquidador chileno para recuperar fondos desviados \\
\addlinespace
Astaldi & 11.\textsuperscript{o} Juzgado Civil de Santiago (2020)
& Reconocimiento de la reorganización de la matriz italiana y coordinación con las
propuestas locales sometidas a votación
& Criterio de convergencia procesal para evitar el desmedro de los acreedores locales \\
\addlinespace
LATAM Airlines Group & 2.\textsuperscript{o} Juzgado Civil de Santiago (2020)
& Reconocimiento del procedimiento del Capítulo 11 estadounidense como procedimiento
extranjero principal, con control sustantivo sobre los bienes locales
& Primer reconocimiento con controversias complejas de deuda y orden de informar el
inventario de bienes en el país \\
\bottomrule
\end{tabularx}
\end{table}

\verificar{Confirmar tribunal, rol, fechas y estado procesal de cada caso antes de
publicar.}

%==============================================================================
\bloqueJurisprudencial
%==============================================================================

\subsection{Límites de las medidas de amparo locales}

Doctrina conforme a la cual el reconocimiento de un procedimiento extranjero principal no
faculta al juez chileno para conceder automáticamente la \pfc{} general del
\art{57 \Nro 1}, debiendo sujetarse al estatuto especial de medidas y auxilios
provisionales, e impidiendo interpretaciones extensivas que vulneren los privilegios de
los acreedores chilenos disidentes.\verificar{Confirmar el artículo aplicable y la fuente
de la doctrina invocada.}
